\section{范数球几何下的推广}\label{sec:norm-geometry}
前两节的端点系数来自径向集中与共同核心。下面先验证所需统一弱型界，
再将这些步骤推广至任意范数球。
令 $K$ 为 $\mathbb R^n$ 中的范数单位球，$\rho(x)=\|x\|_K$，
$K(x,r)=x+rK$。
\begin{lemma}\label{lem:norm-weak}
其中心极大算子对每个非零局部有限 Radon 测度 $\nu$ 满足下式，
其中可取 $B=3^n$：
\[\nu\{M_{K,\nu}f>\alpha\}\le\frac B\alpha\int|f|\,d\nu,
\qquad f\in L^1(\nu),\quad\alpha>0.\tag{N.1}\]
平均中零质量球仍予忽略。
\end{lemma}
\begin{proof}
先估计 Besicovitch 相交常数。对有限个共同相交且每个闭球不含其他球中心的
球族，将公共点平移至零。单球情形显然；其余情形各中心 $x_i\ne0$。
令 $a_i=\rho(x_i)\le r_i$、$u_i=x_i/a_i$。若 $a_i\ge a_j$ 且
$\rho(u_i-u_j)\le1$，则
\[\rho(x_i-x_j)\le(a_i-a_j)\rho(u_i)+a_j\rho(u_i-u_j)\le a_i\le r_i,\]
与球族条件矛盾。因此 $\rho(u_i-u_j)>1$；集合
$u_i+\frac12\operatorname{int}K$ 两两不交且包含于 $\frac32K$。
体积比较给球数至多 $3^n$，对任意球族也可通过有限子族得到相同交叠界。
由 Aldaz 的统一弱型定理 \cite[Theorem 1.2]{Aldaz2020}，该相交常数控制全部
局部有限测度的中心弱型常数，故得 (N.1)。有限维范数空间可分，所引定理的
测度条件在此满足。这里使用的是相交常数与弱型范数的联系，没有把它当作覆盖选择常数。
\end{proof}
令 $M_K=M_{K,dx}$，
\[d\mu_{K,\tau}=\rho^{\tau-n}dx,\quad
\kappa=\frac{n|K|}{\tau},\quad 0<\tau<n,\]
并用 $M_K$ 的严格光晕及 $\mu_{K,\tau}$ 定义外层常数
$C^{\rm out}_{K,\tau}$，输入仍同时满足有限正 Lebesgue 体积及有限正加权质量。
\begin{theorem}\label{thm:norm-power}
对任意范数单位球 $K$，
\[C^{\rm out}_{K,\tau}(1-t)-1\sim t^{\tau/n}.\tag{N.2}\]
若进一步 $0<\tau<T_n$，其中 $T_n$ 如定理 \ref{thm:power-second}，则
\[C^{\rm out}_{K,\tau}(1-t)-1
=t^{\tau/n}+t^{2\tau/n}+o_{K,\tau}(t^{2\tau/n}).\tag{N.3}\]
\end{theorem}
\begin{proof}
先由 (N.1) 的测度扰动得到
\[|\{M_K\1_F>1-t\}\setminus F|\le B\frac t{1-t}|F|.\tag{N.4}\]
具体地，对光晕中与 $F$ 不交的紧集 $J$，令
$d\nu_c=\1_Fdx+c\1_Jdx$。严格见证给
$M_{K,\nu_c}\1_F>[1+ct/(1-t)]^{-1}$ 于 $J$。
代入 (N.1)，令 $c\to\infty$ 并用内正则性，即得 (N.4)。

径向积分和等体积交换给出
\[\mu_{K,\tau}(rK)=\kappa r^\tau,\qquad
\mu_{K,\tau}(A)\le\kappa\left(\frac{|A|}{|K|}\right)^{\tau/n}.\tag{N.5}\]
固定 $C>2$，按定理 \ref{thm:op-power-n-leading} 的方式，把外部光晕分为
存在 $r<C\rho(x)$ 严格见证的短类 $S$ 与其余长类 $T$。
范数球指标函数的平移及半径变化在 $L^1(dx)$ 中连续，故可测性和排除半径等号的论证相同。
短类的二进环带使用 $D_k=2(1+C)2^kK$。由 (N.4) 及同一几何级数，
\[\frac{\mu_{K,\tau}(S)}{\mu_{K,\tau}(E)}
\le\frac{B[2(1+C)]^{n-\tau}}{1-2^{\tau-n}}\frac t{1-t}.\tag{N.6}\]
令
\[z=t^{\tau/n},\qquad A_C=\inf_{\rho(u)\le C^{-1}}
\frac{\mu_{K,\tau}(u+K)}\kappa.\]
在 $A_C>z$ 时，每个长类实际严格见证先给出
$\mu_{K,\tau}(E)\ge\kappa r^\tau(A_C-z)$，故其半径上确界 $R$ 有限。
三角不等式保证 $K(x,r)\supset(r-\rho(x))K$；
趋近 $R$ 的实际见证因此包含公共核心 $(R/C)K$，而
$T\subset(R/C)K\setminus E$。于是 $|T|\le t|K|R^n$，再用 (N.5)，得
\[\frac{\mu_{K,\tau}(T)}{\mu_{K,\tau}(E)}
\le\frac z{A_C-z}.\tag{N.7}\]
若长类为空，则其质量界取零。这些步骤没有先假设存在最大球。

为估计平移，令 $\delta=C^{-1}\le1/2$。层析表示为
\[\mu_{K,\tau}(A)=(n-\tau)\int_0^\infty|A\cap sK|s^{\tau-n-1}\,ds.\tag{N.8}\]
若 $\rho(u)\le\delta$，交叠差
$|K\cap sK|-|(u+K)\cap sK|$ 非负，在 $[1-\delta,1+\delta]$ 外为零。
壳内两个交集都含 $(1-\delta)K$，故差至多为
$|K|[1-(1-\delta)^n]\le n|K|\delta$。积分并除以 $\kappa$，得
\[0\le1-A_C\le\tau(n-\tau)2^{n+2-\tau}C^{-2}.\tag{N.9}\]
因此 (N.6)--(N.7) 分别给出与 (S.4) 同形的短类误差和长类包络。
取 $0<a<1/n$、$C=t^{-a}$ 得到首项上界；取 (S.5) 的 $a$ 则得到二阶上界。

下界取 $E=RK\setminus hK$，$h/R=t^{1/n}/(1+t^{1/n})$。
对 $\rho(x)<h$ 取 $r=R-\rho(x)$，三角不等式同时保证
$K(x,r)\subset RK$ 且在 $3h<R$ 时包含 $hK$，缺失比例严格小于 $t$。
球面为 Lebesgue 零集，环的内部与孔内均属严格光晕，故下界仍为
$z/[(1+t^{1/n})^\tau-z]$。由 (S.6) 的展开，分别与首项和指定范围的二阶上界匹配。
\end{proof}

\begin{corollary}[共同中心多集合的几何传递]
以同一 $K(x,r)$ 平均各指标函数，并按正文公共交集归一化。
对 $P_{\boldsymbol a}(1)<\lambda<1$，有
\[C^{\boldsymbol a,\rm out}_{K,\tau}(\lambda)
=C^{\rm out}_{K,\tau}(1-s_{\boldsymbol a}(\lambda)).\tag{N.10}\]
因而首项系数为 $a_*^{-\tau/n}$；在 $\tau<T_n$ 时第二项系数为
$a_*^{-2\tau/n}$，其中 $a_*=\min_i a_i$。
\end{corollary}
\begin{proof}
定理 \ref{thm:weighted-transfer} 的正向缺陷估计与几何形状无关。
反向构造中，凸体球面为 Lebesgue 零集，中心移动的球对称差体积趋零，
故有限参考球仍有稳定的严格余量。对这些球生成的有限 Boolean 原子作
Lebesgue 比例分割，便得到原证明所需的共同输入；光晕测量用的
$\mu_{K,\tau}\ll dx$ 不改变严格见证。取有界 $D$ 包住有限参考球及其扰动球，在 $D\setminus F$ 的 Boolean 原子上分割，
再将 $D^c\cap F^c$ 延拓为不交各部 $A_i$，令 $E_i=\mathbb R^n\setminus A_i$。
于是公共交集恰为 $F$；由加权内正则性穷尽严格光晕，
给出 (N.10)。最后代入第一水位展开即得系数。
\end{proof}
此推广使用范数球的径向结构及已验证的 (N.1)，未把一般 Besicovitch 几何族替换为范数球。
不同范数下的权与平均球须同时匹配；本文没有把欧氏径向权代入方体算子得到同一系数。
